\documentclass[11pt]{article}
\usepackage[T1]{fontenc}
\usepackage{lmodern}
\usepackage[margin=1.04in]{geometry}
\usepackage{amsmath,amssymb,amsthm,mathtools}
\usepackage{microtype,xcolor}
\usepackage[hyphens]{url}
\usepackage[colorlinks=true,linkcolor=blue!45!black,citecolor=blue!45!black,urlcolor=blue!45!black]{hyperref}
\newcommand{\E}{\mathbb E}
\newcommand{\R}{\mathbb R}
\DeclareMathOperator{\Var}{Var}
\newtheorem{theorem}{Theorem}[section]
\newtheorem{lemma}[theorem]{Lemma}
\newtheorem{proposition}[theorem]{Proposition}
\newtheorem{corollary}[theorem]{Corollary}
\theoremstyle{remark}
\newtheorem{remark}[theorem]{Remark}
\numberwithin{equation}{section}
\hypersetup{pdftitle={A uniform influence bound for hypergraph properties},pdfauthor={OpenAI},pdfsubject={Variance and pivotal influence for symmetric Boolean hypergraph properties}}
\title{A uniform influence bound for hypergraph properties}
\author{OpenAI}
\date{October 5, 2026}
\begin{document}
\maketitle
\begin{abstract}
For every fixed integer $r\ge3$, we prove that every relabeling-invariant
Boolean property of simple $r$-uniform hypergraphs on $n$ vertices
satisfies
$\Var_p(f)\le C_r I_p(f)/(\log n)^{r/(r-1)}$.
The constant depends only on $r$, and the bound holds uniformly for all
$0<p<1$ without a monotonicity assumption. For increasing properties,
it gives the corresponding threshold-width bound with exponent
$r/(r-1)$, proving the hypergraph threshold-width conjecture of
Friedgut and Kalai.
\end{abstract}
\section{Introduction}
\label{sec:introduction}

Let $r\ge3$ and $n\ge r$ be integers, write $[n]=\{1,\ldots,n\}$,
and put $E_{n,r}=\binom{[n]}r$.
An element of $\{0,1\}^{E_{n,r}}$ represents a simple $r$-uniform
hypergraph. A \emph{hypergraph property} is a Boolean function on this
space that is invariant under every permutation of $[n]$.
For $0<p<1$, give the coordinates the independent Bernoulli$(p)$ law,
and denote expectation and variance by $\E_p$ and $\Var_p$.

For an edge $e$, let $x_{e\leftarrow b}$ be $x$ with its $e$-coordinate
set to $b$. Define the \emph{pivotal influence} of $e$ and the
\emph{total influence} of $f$ by
\[
 I_{e,p}(f)=\Pr_p\bigl(f(x_{e\leftarrow1})\ne
                         f(x_{e\leftarrow0})\bigr),
 \qquad I_p(f)=\sum_{e\in E_{n,r}}I_{e,p}(f).
\]
The probability in the first expression is over the other coordinates;
it does not include an additional factor of $p(1-p)$.
All logarithms are natural, and $n$ always counts vertices.

\begin{theorem}\label{thm:main}
For every integer $r\ge3$ there is a finite constant $C_r$ such that,
for every integer $n\ge r$, every $0<p<1$, and every hypergraph property
$f:\{0,1\}^{E_{n,r}}\to\{0,1\}$,
\begin{equation}\label{eq:main}
 \Var_p(f)\le \frac{C_r}{(\log n)^{r/(r-1)}}I_p(f).
\end{equation}
No monotonicity assumption is required.
\end{theorem}

The theorem measures how vertex symmetry forces sensitivity to edge
changes. If the probability of a property stays bounded away from
zero and one, its total influence must grow at least as
$(\log n)^{r/(r-1)}$.

A property is \emph{increasing} if changing coordinates from $0$ to
$1$ cannot decrease its value, and \emph{nontrivial} if it is not
constant. For a nontrivial increasing property, put $q(p)=\E_p f$
and define $p_a=\inf\{p\in[0,1]:q(p)\ge a\}$ for $0<a<1$.
For fixed $0<\varepsilon<1/2$, its \emph{threshold width} is
$p_{1-\varepsilon}-p_\varepsilon$.
Friedgut and Kalai conjectured that this width is at most
$O_{r,\varepsilon}((\log n)^{-r/(r-1)})$ for every increasing
$r$-uniform hypergraph property~\cite[Section~5, p.~2999]{FK96}.
Their formulation counts edge coordinates rather than vertices;
for fixed $r$, $\log\binom nr\sim r\log n$ as $n\to\infty$.
Corollary~\ref{cor:threshold} proves this conjectured bound by
integrating~\eqref{eq:main}, with an explicit dependence on
$\varepsilon$.

\paragraph{Context and attribution.}
Kahn, Kalai, and Linial~\cite{KKL88} established a fundamental link
between Fourier analysis and the influences of Boolean functions.
For a function invariant under a transitive group on $M$ coordinates,
their theorem gives a total-influence lower bound of order
$\Var_{1/2}(f)\log M$ at the uniform measure.
Bourgain, Kahn, Kalai, Katznelson, and Linial~\cite{BKKKL92}
extended the influence method to product spaces.
Friedgut and Kalai~\cite{FK96} used influence inequalities to prove
sharp thresholds for increasing families with transitive symmetry.
For hypergraph properties, the induced action on the
$\binom nr$ edge coordinates is transitive. The richer structure of
vertex permutations, however, contains more information than this
single transitivity statement.

Bourgain and Kalai~\cite[Corollary~4.1]{BK97} studied how group
actions on sets of coordinates improve influence estimates.
For fixed $r$ and every $\eta>0$, their result gives a lower bound
of order $\Var_{1/2}(f)(\log n)^{r/(r-1)-\eta}$ for the total
influence of a hypergraph property at the uniform measure.
Their influence results already allow nonmonotone functions.
Kelman, Kindler, Lifshitz, Minzer, and Safra obtained further
estimates for group-invariant functions~\cite{KKLMS20}.
For graph properties, these give
$I_{1/2}(f)\ge c\Var_{1/2}(f)(\log n)^2/(\log\log n)^2$
for sufficiently large $n$; see the arXiv version cited
in~\cite[Section~7.4]{KKLMS20}.
The present proof uses vertex symmetry through restrictions: all
edges touching a vertex block remain free, and only edges wholly
outside the block are fixed. This is the method of the companion
manuscript \emph{A Sharp Threshold Bound for Monotone Graph
Properties}~\cite[Theorem~1.2 and Sections~2--3]{Graph2026}, which establishes
the corresponding variance bound with denominator $(\log n)^2$
for graph properties, uniformly in $p$ and without monotonicity.
We reproduce its Fourier estimate and restriction argument in full.
The additional combinatorial step replaces the graph degree estimate
by its $r$-uniform counterpart, producing the exponent $r/(r-1)$.

\paragraph{Proof strategy.}
Put $\sigma=\sqrt{p(1-p)}$ and $\alpha=r/(r-1)$.
The Fourier expansion for the Bernoulli product measure is indexed by
sets $S$ of hyperedges. Its squared nonconstant coefficients sum to
the variance; weighting them by $|S|$ gives $\sigma^2I_p(f)$.
Thus the influence controls high degrees. The task is to bound the
remaining mass using symmetry.

For large $n$, choose a uniformly random block $B$ of $m$ vertices,
with $m$ about $\sqrt n$, and leave free all edges meeting $B$.
Every permutation within $B$ survives an assignment to the other
coordinates, and every free-edge orbit has at least $m$ elements.
Section~\ref{sec:fourier} bounds the restricted Fourier mass up to
a degree $k$ proportional to $\sigma\log m$ by $m^{-1/4}$ times
the restricted total influence. Its linear dependence on influence
allows averaging over the assigned coordinates.

Regard a nonempty Fourier index $S$ as an $r$-uniform hypergraph
with $s$ edges. It has a vertex of positive degree at most
$r s^{(r-1)/r}$. For $s\le L=(k/r)^\alpha$, this degree is at
most $k$. The event that $B$ meets the union of these edges only
at that vertex has probability comparable to $m/n$ and leaves
between one and $k$ support edges free. After averaging over the
assignments and $B$, the restricted total influence is at most
$(rm/n)I_p(f)$, so the two factors of $m/n$ cancel.
Section~\ref{sec:restriction} thereby controls the original
Fourier mass up to degree $L$.

Since $L$ is of order $(\sigma\log n)^\alpha$, the remaining
mass is at most $\sigma^2I_p(f)/L$, with coefficient of order
$\sigma^{2-\alpha}(\log n)^{-\alpha}$. The positive exponent
$2-\alpha$ makes this bound uniform in $p$.
Section~\ref{sec:completion} completes the estimate and derives
the threshold consequence.

\section{A Fourier estimate for small individual influences}
\label{sec:fourier}

We first work on an arbitrary finite coordinate set $J$. Our objective
is to control low-degree Fourier mass when every individual influence
is small relative to the total. We follow the elementary argument
in~\cite[Lemma~2.1]{Graph2026}; its analytic ingredient is a
two-to-four norm estimate of the type underlying the classical
hypercontractive inequalities~\cite{Bonami70,Beckner75}.
The proof below includes the dependence on the bias.

Fix $0<p<1$ and put $\sigma=\sqrt{p(1-p)}$. For $i\in J$ and
$A\subseteq J$, define
\[
 \chi_i(x)=\frac{x_i-p}{\sigma},\qquad
 \chi_A(x)=\prod_{i\in A}\chi_i(x),\qquad \chi_\varnothing=1.
\]
Independence, zero means, and unit second moments show that the
$\chi_A$ form an orthonormal basis. Write
$\widehat h(A)=\E_p[h\chi_A]$.
Throughout this section, norms and expectations use this product law,
and $I_i(h)$ and $I(h)$ denote pivotal and total influence on $J$.

For Boolean $h$, let $\Delta_i h$ be the function on the other
coordinates obtained by subtracting the section at $x_i=0$ from the
section at $x_i=1$. Direct expectation in coordinate $i$ gives
\[
 \widehat h(A\cup\{i\})=
 \sigma\widehat{\Delta_i h}(A)\qquad(A\subseteq J\setminus\{i\}).
\]
Because $\Delta_i h\in\{-1,0,1\}$, Parseval's identity yields
\begin{equation}\label{eq:parseval}
 \sum_{S\subseteq J}|S|\widehat h(S)^2=\sigma^2 I(h),
 \qquad
 \Var_p(h)=\sum_{\varnothing\ne S\subseteq J}\widehat h(S)^2.
\end{equation}
In particular, this identity uses no sign condition on $\Delta_i h$.

\begin{lemma}[Low-degree mass]\label{lem:low}
Let $m\ge2$ be real and let $h:\{0,1\}^J\to\{0,1\}$ satisfy
$I_i(h)\le I(h)/m$ for every $i\in J$. With
$k=\sigma\log m/64$, one has
\begin{equation}\label{eq:low}
 \sum_{1\le |S|\le k}\widehat h(S)^2\le m^{-1/4}I(h).
\end{equation}
The sum is empty when $k<1$.
\end{lemma}

\begin{proof}
We will use a two-to-four norm estimate to bound low-degree projections
of the differences $\Delta_i h$.
Set $\rho=\sigma/4$ and define the linear operator $T$ by
$T\chi_A=\rho^{|A|}\chi_A$. We first prove
\begin{equation}\label{eq:24}
 \|Tg\|_4\le\|g\|_2
\end{equation}
for every real-valued function $g$ on the product space.
For one coordinate, $|\chi_i|\le1/\sigma$ and $\E\chi_i^2=1$
imply $|\E\chi_i^3|\le1/\sigma$ and
$\E\chi_i^4\le1/\sigma^2$. For real $a,b$, expansion and
$2|ab^3|\le a^2b^2+b^4$ give
\begin{align*}
 \E(a+\rho b\chi_i)^4
 &\le a^4+6\rho^2a^2b^2+
       \frac{2\rho^3}{\sigma}(a^2b^2+b^4)
       +\frac{\rho^4}{\sigma^2}b^4\\
 &=a^4+\frac{13\sigma^2}{32}a^2b^2
       +\frac{9\sigma^2}{256}b^4
 \le(a^2+b^2)^2,
\end{align*}
where the last inequality uses $\sigma\le1/2$.

To extend this estimate to the product, induct on $|J|$, starting with
the constant functions when $J$ is empty. Write $g=a+b\chi_i$,
where $a,b$ depend on the remaining coordinates, and use $T$ for the
operator on those coordinates as well. The one-coordinate estimate,
applied pointwise, followed by the $L^2$ triangle inequality gives
\[
 \|Tg\|_4^2\le\|(Ta)^2+(Tb)^2\|_2
 \le\|Ta\|_4^2+\|Tb\|_4^2
 \le\|a\|_2^2+\|b\|_2^2=\|g\|_2^2.
\]
This proves~\eqref{eq:24}.

Let $P_{\le k}$ be the orthogonal projection onto the span of
$\chi_A$ with $|A|\le k$. For every $u$ in this span,
\[
 \|u\|_4\le\|T^{-1}u\|_2\le\rho^{-k}\|u\|_2.
\]
Taking the supremum of $|\E[gu]|$ over unit $L^2$ vectors in this
span, and applying H\"older's inequality, therefore gives
\begin{equation}\label{eq:projection}
 \|P_{\le k}g\|_2\le\rho^{-k}\|g\|_{4/3}.
\end{equation}
These inequalities are valid for real $k\ge0$, with no rounding
required.

Each nonempty Fourier support is counted at least once when summed
over its coordinates. Using the difference identity and enlarging
the resulting degree cutoff from $k-1$ to $k$, we obtain
\begin{align}
 \sum_{1\le|S|\le k}\widehat h(S)^2
 &\le\sigma^2\sum_{i\in J}
       \|P_{\le k}\Delta_i h\|_2^2 \notag\\
 &\le\sigma^2\rho^{-2k}\sum_{i\in J}I_i(h)^{3/2}.
 \label{eq:derivative-mass}
\end{align}
Here each projection acts on the coordinates other than $i$, and
$\|\Delta_i h\|_{4/3}^2=I_i(h)^{3/2}$ because
$|\Delta_i h|$ is an indicator.

The function $\sigma\log(4/\sigma)$ is increasing on $(0,1/2]$
and is at most $2$. Hence $\rho^{-2k}\le m^{1/16}$.
If $I(h)\le m^{1/4}$, then
$I_i(h)^{1/2}\le m^{-3/8}$, so~\eqref{eq:derivative-mass}
is at most $\sigma^2m^{-5/16}I(h)\le m^{-1/4}I(h)$.
If $I(h)>m^{1/4}$, the right side of~\eqref{eq:low} exceeds $1$,
whereas its left side is at most $\Var_p(h)\le1/4$.
\end{proof}

\section{Restriction to a vertex block}
\label{sec:restriction}

Return to a hypergraph property $f$ on $E=E_{n,r}$ and fix the bias
$p$, with $\sigma=\sqrt{p(1-p)}$. Its Fourier coefficients use the basis of
Section~\ref{sec:fourier}, with $J=E$.
The restriction method of~\cite[Lemma~3.1 and Proposition~3.2]{Graph2026} preserves a
large permutation group after conditioning. We first obtain an
averaged estimate for restricted Fourier degrees, and then relate
those degrees to the sizes of the original supports.

For a vertex set $B\subseteq[n]$, write
\[
 F_B=\{e\in E:e\cap B\ne\varnothing\}
\]
for the edge coordinates that remain free.

\begin{lemma}[Averaged restriction estimate]\label{lem:block}
Let $r+1\le m\le n$ be integers and set $k=\sigma\log m/64$.
For $B$ chosen uniformly among the $m$-subsets of $[n]$,
\begin{equation}\label{eq:block}
 \sum_{S\subseteq E}
 \Pr_B(1\le|S\cap F_B|\le k)\widehat f(S)^2
 \le m^{-1/4}\frac{rm}{n}I_p(f).
\end{equation}
\end{lemma}

\begin{proof}
Fix $B$ and an assignment $y$ to the coordinates in $E\setminus F_B$.
Denote the resulting Boolean function on $F_B$ by $f_y$.
Every permutation supported on $B$ fixes each edge outside $F_B$
individually, and consequently fixes $y$. Thus $f_y$ is invariant
under all these permutations, regardless of the assignment.

An edge with $j$ vertices in $B$ has an orbit of size $\binom mj$:
its outside part is fixed and its inside part can be any $j$-subset
of $B$. Since $1\le j\le r<m$, this size is at least $m$.
The free product measure is invariant under the same permutations,
so influences are equal within each orbit. It follows that
\[
 I_e(f_y)\le I(f_y)/m\qquad(e\in F_B).
\]
Lemma~\ref{lem:low} therefore applies to every $f_y$.

Average its conclusion over the product law of $y$. For $A\subseteq F_B$,
the Fourier expansion of $f$ gives
\[
 \widehat{f_y}(A)=
 \sum_{D\subseteq E\setminus F_B}\widehat f(A\cup D)\chi_D(y).
\]
Orthonormality in the assigned coordinates yields
\[
 \E_y\bigl[\widehat{f_y}(A)^2\bigr]
 =\sum_{D\subseteq E\setminus F_B}\widehat f(A\cup D)^2.
\]
Also $\E_y I_e(f_y)=I_{e,p}(f)$ for every $e\in F_B$, by
conditioning the disagreement event. Consequently
\[
 \sum_{S:\,1\le|S\cap F_B|\le k}\widehat f(S)^2
 \le m^{-1/4}\sum_{e\in F_B}I_{e,p}(f).
\]
Finally average over $B$. For each fixed edge $e$, a union bound on
its $r$ vertices gives $\Pr_B(e\in F_B)\le rm/n$, proving
\eqref{eq:block}.
\end{proof}

The next lemma extends the graph capture argument
of~\cite[Lemma~3.3]{Graph2026} to identify the hypergraph supports
detected by this estimate. Here a support $S$ is itself a set of distinct $r$-edges;
its vertex union contains no isolated vertices.

\begin{lemma}[Capturing a support at one vertex]\label{lem:capture}
Let $n\ge r\ge3$ and $1\le m\le n$ be integers, let $k>0$, and put
\[
 \alpha=\frac r{r-1},\qquad L=(k/r)^\alpha.
\]
Suppose
\begin{equation}\label{eq:capture-condition}
 rL\frac{m-1}{n-1}\le\frac12.
\end{equation}
For a uniform $m$-subset $B$ of $[n]$ and every
$S\subseteq E_{n,r}$ with $1\le|S|\le L$,
\begin{equation}\label{eq:capture}
 \Pr_B(1\le|S\cap F_B|\le k)\ge\frac{m}{2n}.
\end{equation}
\end{lemma}

\begin{proof}
Fix such an $S$, write $s=|S|$, and let $V=\bigcup_{e\in S}e$
and $v=|V|$. Because $s\le\binom vr\le v^r$, the average degree
of the hypergraph with edge set $S$ is at most
\[
 \frac{rs}{v}\le r s^{(r-1)/r}\le k.
\]
Choose a vertex $u\in V$ of degree $d_S(u)$ between $1$ and $k$.
This choice depends only on $S$, before $B$ is sampled.

Conditional on $u\in B$, each other vertex lies in $B$ with
probability $(m-1)/(n-1)$. The union bound and $v\le rs\le rL$
therefore give
\begin{align*}
 \Pr_B(B\cap V=\{u\})
 &\ge\frac mn\left(1-(v-1)\frac{m-1}{n-1}\right)\\
 &\ge\frac m{2n}.
\end{align*}
On this event, the edges of $S$ meeting $B$ are exactly those
containing $u$. Thus $|S\cap F_B|=d_S(u)\in[1,k]$.
\end{proof}

The captured set need not consist of a single edge: the block meets
one support vertex, and retains every support edge incident with it.
Also the vertex $u$ may vary with $S$, since~\eqref{eq:capture}
is a separate probability bound for each nonnegative summand
in~\eqref{eq:block}.

Combining Lemmas~\ref{lem:block} and~\ref{lem:capture} gives,
whenever their hypotheses hold,
\begin{equation}\label{eq:captured-mass}
 \sum_{1\le|S|\le L}\widehat f(S)^2
 \le 2r m^{-1/4}I_p(f).
\end{equation}
If $L<1$, this conclusion holds because the sum is empty.
Otherwise, restrict the nonnegative sum in~\eqref{eq:block} to
$1\le|S|\le L$ and divide by the lower bound $m/(2n)$.
This completes the transfer from restricted degree $k$ to original
degree $(k/r)^{r/(r-1)}$.

\section{The uniform bound and its threshold consequence}
\label{sec:completion}

We now choose the block size so that~\eqref{eq:captured-mass}
controls the low-degree part of the variance. The weighted Parseval
identity will control the remaining degrees.

\begin{proof}[Proof of Theorem~\ref{thm:main}]
Fix $r\ge3$ and write $\alpha=r/(r-1)$. Choose an integer $N_r>r$
such that, for every $n\ge N_r$ and $m=\lfloor\sqrt n\rfloor$,
\begin{equation}\label{eq:large-n}
 \begin{gathered}
 m\ge r+1,\qquad \log m\ge\tfrac13\log n,\\
 r(\log n)^\alpha n^{-1/2}\le\tfrac12,\qquad
 m^{-1/4}(\log n)^\alpha\le1.
 \end{gathered}
\end{equation}
Such an $N_r$ exists: $m$ is asymptotic to $\sqrt n$, and any
fixed power of $\log n$ grows more slowly than a positive power of
$n$. In particular, the choice is independent of $p$ and $f$.

Assume first that $n\ge N_r$, and set
\[
 \sigma=\sqrt{p(1-p)},\qquad
 k=\frac{\sigma\log m}{64},\qquad L=(k/r)^\alpha.
\]
Both cutoffs are positive. Since $k/r\le\log n$ and
\[
 \frac{m-1}{n-1}\le\frac{\sqrt n-1}{n-1}
 =\frac1{\sqrt n+1}\le n^{-1/2},
\]
the third condition in~\eqref{eq:large-n} implies
\eqref{eq:capture-condition}. Thus~\eqref{eq:captured-mass} and
the fourth condition in~\eqref{eq:large-n} give
\begin{equation}\label{eq:low-final}
 \sum_{1\le|S|\le L}\widehat f(S)^2
 \le\frac{2r}{(\log n)^\alpha}I_p(f).
\end{equation}

For the remaining degrees,~\eqref{eq:parseval} gives
\[
 \sum_{|S|>L}\widehat f(S)^2
 \le\frac1L\sum_S|S|\widehat f(S)^2
 =\frac{\sigma^2}{L}I_p(f).
\]
The decisive bias dependence is
\begin{equation}\label{eq:bias-cancellation}
 \frac{\sigma^2}{L}
 =\frac{(64r)^\alpha\sigma^{2-\alpha}}{(\log m)^\alpha}
 \le\frac{(192r)^\alpha}{(\log n)^\alpha}.
\end{equation}
Indeed $2-\alpha=(r-2)/(r-1)>0$, $\sigma\le1$, and
$\log m\ge(\log n)/3$. This estimate remains valid at arbitrarily
small positive $\sigma$. If $L<1$, the low-degree sum is empty and
the same tail estimate bounds the whole variance.

Adding the two parts proves~\eqref{eq:main} for $n\ge N_r$ with
constant $2r+(192r)^\alpha$. For $r\le n<N_r$,
\eqref{eq:parseval} gives $\Var_p(f)\le\sigma^2 I_p(f)\le I_p(f)/4$.
Therefore a valid choice for all $n$ and $p$ is
\[
 C_r=\max\left\{2r+(192r)^\alpha,\,
                       \frac{(\log N_r)^\alpha}{4}\right\}.
\]
\end{proof}

We finish with the threshold consequence stated in the introduction.
Recall that $q(p)=\E_p f$ and
$p_a=\inf\{p\in[0,1]:q(p)\ge a\}$.

\begin{corollary}\label{cor:threshold}
For a nontrivial increasing property of $r$-uniform hypergraphs on
$n\ge r$ vertices, where $r\ge3$, and every $0<\varepsilon<1/2$,
\[
 p_{1-\varepsilon}-p_\varepsilon
 \le\frac{2C_r}{(\log n)^{r/(r-1)}}
                  \log\frac{1-\varepsilon}{\varepsilon}.
\]
\end{corollary}

\begin{proof}
The Margulis--Russo identity relates pivotal influences to the
probability derivative; see~\cite{Margulis74},
\cite[Lemma~3]{Russo81}, and \cite[Lemma~2.2]{FK96}.
Here it follows directly by giving each coordinate
its own bias: the expectation is affine in each bias separately,
and its partial derivative in coordinate $e$ is
$\E_p[f(x_{e\leftarrow1})-f(x_{e\leftarrow0})]$.
Monotonicity identifies this with $I_{e,p}(f)$. Differentiating
along the common bias gives $q'(p)=I_p(f)$.

Nontrivial monotonicity gives $q(0)=0$ and $q(1)=1$. Some pair of
configurations differing in one coordinate has different values,
so its pivotal event has positive probability for every $0<p<1$.
Hence $q'(p)>0$ there, and the stated quantiles exist uniquely.
Theorem~\ref{thm:main} implies
\[
 \frac{q'(p)}{q(p)(1-q(p))}
 \ge\frac{(\log n)^{r/(r-1)}}{C_r}.
\]
Integrating from $p_\varepsilon$ to $p_{1-\varepsilon}$, using
the antiderivative $\log(q/(1-q))$, proves the claim.
\end{proof}

\clearpage
\begin{thebibliography}{10}

\bibitem{Beckner75}
W. Beckner, Inequalities in Fourier analysis,
\emph{Annals of Mathematics} \textbf{102} (1975), 159--182.

\bibitem{Bonami70}
A. Bonami, \'{E}tude des coefficients de Fourier des fonctions de
$L^p(G)$, \emph{Annales de l'Institut Fourier} \textbf{20} (1970),
no.~2, 335--402.

\bibitem{BKKKL92}
J. Bourgain, J. Kahn, G. Kalai, Y. Katznelson, and N. Linial,
The influence of variables in product spaces,
\emph{Israel Journal of Mathematics} \textbf{77} (1992), 55--64.

\bibitem{BK97}
J. Bourgain and G. Kalai, Influences of variables and threshold
intervals under group symmetries,
\emph{Geometric and Functional Analysis} \textbf{7} (1997), 438--461.

\bibitem{FK96}
E. Friedgut and G. Kalai, Every monotone graph property has a sharp
threshold, \emph{Proceedings of the American Mathematical Society}
\textbf{124} (1996), 2993--3002.

\bibitem{KKL88}
J. Kahn, G. Kalai, and N. Linial, The influence of variables on Boolean
functions, in \emph{Proceedings of the 29th Annual Symposium on
Foundations of Computer Science}, IEEE, 1988, 68--80.

\bibitem{KKLMS20}
E. Kelman, G. Kindler, N. Lifshitz, D. Minzer, and M. Safra,
Towards a proof of the Fourier-Entropy Conjecture?,
\emph{Geometric and Functional Analysis} \textbf{30} (2020), 1097--1138.
See also arXiv:1911.10579v2, May 7, 2020.

\bibitem{Margulis74}
G. A. Margulis, Probabilistic characteristics of graphs with large
connectivity, \emph{Problems of Information Transmission}
\textbf{10} (1974), 174--179.

\bibitem{Graph2026}
OpenAI, \emph{A Sharp Threshold Bound for Monotone Graph Properties},
OpenAI Math Release preprint
\href{https://github.com/openai/math/blob/main/preprints/A-Sharp-Threshold-Bound-for-Monotone-Graph-Properties-September-25-2026/paper.pdf}{OAI:A-Sharp-Threshold-Bound-for-Monotone-Graph-Properties-September-25-2026},
2026.

\bibitem{Russo81}
L. Russo, On the critical percolation probabilities,
\emph{Zeitschrift f\"ur Wahrscheinlichkeitstheorie und verwandte Gebiete}
\textbf{56} (1981), 229--237.

\end{thebibliography}

\end{document}
